% --- Archon physics formalization source begin ---
% archon:physics
% archon:covers PhyXMiniProblems/problem_phyx_mini_0273.lean
% archon:source-report reports/phyx_mini/problem_phyx_mini_0273.source.json

\chapter{Physics problem phyx\_mini\_0273}
\label{ch:PhyXMiniProblems_problem_phyx_mini_0273}

\paragraph{Problem source.}
\#\# Physical scenario

Figure gives the position \$x(t)\$ of a block oscillating in SHM on the end of a spring (\$t\_s = 40.0\textbackslash{} ms\$).

\#\# Question

What is the magnitude of the radial acceleration of a particle in the corresponding uniform circular motion?

\#\# Answer choices

- A: A: \$1.3 \textbackslash{}times 10\textasciicircum{}3\textbackslash{}\textbackslash{} m/s\textasciicircum{}2\$

- B: B: \$1.5 \textbackslash{}times 10\textasciicircum{}3\textbackslash{}\textbackslash{} m/s\textasciicircum{}2\$

- C: C: \$1.7 \textbackslash{}times 10\textasciicircum{}3\textbackslash{}\textbackslash{} m/s\textasciicircum{}2\$

- D: D: \$1.9 \textbackslash{}times 10\textasciicircum{}3\textbackslash{}\textbackslash{} m/s\textasciicircum{}2\$

\#\# Figure caption (auxiliary; use the image as primary evidence)

The image is a graph depicting a sinusoidal wave. Here are its detailed components:

- The horizontal axis is labeled "t (ms)" representing time in milliseconds.
- The vertical axis is labeled "x (cm)" representing displacement in centimeters.
- The graph shows a sine wave starting at (0,0) and peaks at 8 cm, troughs at -8 cm.
- There are two complete cycles visible.
- A specific point on the horizontal axis is marked as \textbackslash{}( t\_s \textbackslash{}).
- The graph is set on a grid, indicating measurements with smaller increments.
- The axes include tick marks with numerical labels for reference, such as 8, 4, 0, -4, and -8 on the y-axis.

\#\# Dataset metadata

Wave Properties; Physical Model Grounding Reasoning, Multi-Formula Reasoning

\paragraph{Recorded answer/context.}
C: C: \$1.7 \textbackslash{}times 10\textasciicircum{}3\textbackslash{}\textbackslash{} m/s\textasciicircum{}2\$

\paragraph{Figure/image path.}
/root/proposal\_for\_physic/science-mango/phyx\_mini\_run/phyx\_data/test\_image/273.png

\paragraph{Formalization target.}
create a compiling Lean file with sorry bodies at `PhyXMiniProblems/problem\_phyx\_mini\_0273.lean`. The Lean declarations must preserve the physical quantities, dimensions or dimensional roles, figure labels, governing-law hypotheses, and final relation expressed by this problem.
Use Mathlib/Physlib names found through LeanExplore where available. If a physics API is missing, introduce faithful local abstractions rather than scalar placeholder aliases.

\begin{theorem}[Physics formalization target]
\label{thm:physics:phyx_mini_0273:target}
\lean{PhyXMiniProblems.ProblemPhyXMini0273.radialAcceleration_exact_and_matches_choiceC}
\uses{def:physics:phyx-mini-0273:phyxminiproblems-problemphyxmini0273-accelerationinmeterspersecondsquared, def:physics:phyx-mini-0273:phyxminiproblems-problemphyxmini0273-shmcircularmotionsetup, def:physics:phyx-mini-0273:phyxminiproblems-problemphyxmini0273-matchesproblemandprimaryfigure, def:physics:phyx-mini-0273:phyxminiproblems-problemphyxmini0273-hasphysicalparameters, def:physics:phyx-mini-0273:phyxminiproblems-problemphyxmini0273-satisfiessimpleharmonicmotionlaws, def:physics:phyx-mini-0273:phyxminiproblems-problemphyxmini0273-satisfiesshmuniformcircularmotioncorrespondence, def:physics:phyx-mini-0273:phyxminiproblems-problemphyxmini0273-satisfiesuniformcircularmotionradialaccelerationlaw, def:physics:phyx-mini-0273:phyxminiproblems-problemphyxmini0273-matchesdisplayedanswer}
The assigned autoformalize agent should translate this physics problem into Lean declarations in the covered file, with theorem and lemma proof bodies written as `by sorry`.
\end{theorem}
\begin{proof}
This is an autoformalization task, not a proof task. Produce statements that can later be proved by the physics prover without weakening the physical model.
\end{proof}

% --- Archon declaration topology begin ---
\section{Declaration topology}
\begin{definition}[Length Magnitude Quantity]
  \label{def:physics:phyx-mini-0273:phyxminiproblems-problemphyxmini0273-lengthmagnitudequantity}
  \lean{PhyXMiniProblems.ProblemPhyXMini0273.LengthMagnitudeQuantity}
  A nonnegative physical length, used for amplitudes and circle radii
\end{definition}
\begin{proof}
  This is the declared model definition of Length Magnitude Quantity.
\end{proof}
\begin{definition}[Signed Length Quantity]
  \label{def:physics:phyx-mini-0273:phyxminiproblems-problemphyxmini0273-signedlengthquantity}
  \lean{PhyXMiniProblems.ProblemPhyXMini0273.SignedLengthQuantity}
  A signed physical length, used for the one-dimensional position x(t)
\end{definition}
\begin{proof}
  This is the declared model definition of Signed Length Quantity.
\end{proof}
\begin{definition}[Time Quantity]
  \label{def:physics:phyx-mini-0273:phyxminiproblems-problemphyxmini0273-timequantity}
  \lean{PhyXMiniProblems.ProblemPhyXMini0273.TimeQuantity}
  A nonnegative physical duration
\end{definition}
\begin{proof}
  This is the declared model definition of Time Quantity.
\end{proof}
\begin{definition}[Angular Frequency Magnitude Quantity]
  \label{def:physics:phyx-mini-0273:phyxminiproblems-problemphyxmini0273-angularfrequencymagnitudequantity}
  \lean{PhyXMiniProblems.ProblemPhyXMini0273.AngularFrequencyMagnitudeQuantity}
  A nonnegative angular frequency; radians are dimensionless
\end{definition}
\begin{proof}
  This is the declared model definition of Angular Frequency Magnitude Quantity.
\end{proof}
\begin{definition}[Acceleration Magnitude Quantity]
  \label{def:physics:phyx-mini-0273:phyxminiproblems-problemphyxmini0273-accelerationmagnitudequantity}
  \lean{PhyXMiniProblems.ProblemPhyXMini0273.AccelerationMagnitudeQuantity}
  A nonnegative acceleration magnitude, with dimension length per time squared
\end{definition}
\begin{proof}
  This is the declared model definition of Acceleration Magnitude Quantity.
\end{proof}
\begin{definition}[length Magnitude Readout]
  \label{def:physics:phyx-mini-0273:phyxminiproblems-problemphyxmini0273-lengthmagnitudereadout}
  \lean{PhyXMiniProblems.ProblemPhyXMini0273.lengthMagnitudeReadout}
  \uses{def:physics:phyx-mini-0273:phyxminiproblems-problemphyxmini0273-lengthmagnitudequantity}
  Read a nonnegative length in a selected length unit
\end{definition}
\begin{proof}
  This is the declared model definition of length Magnitude Readout.
\end{proof}
\begin{definition}[signed Length Readout]
  \label{def:physics:phyx-mini-0273:phyxminiproblems-problemphyxmini0273-signedlengthreadout}
  \lean{PhyXMiniProblems.ProblemPhyXMini0273.signedLengthReadout}
  \uses{def:physics:phyx-mini-0273:phyxminiproblems-problemphyxmini0273-signedlengthquantity}
  Read a signed displacement in a selected length unit
\end{definition}
\begin{proof}
  This is the declared model definition of signed Length Readout.
\end{proof}
\begin{definition}[time Readout]
  \label{def:physics:phyx-mini-0273:phyxminiproblems-problemphyxmini0273-timereadout}
  \lean{PhyXMiniProblems.ProblemPhyXMini0273.timeReadout}
  \uses{def:physics:phyx-mini-0273:phyxminiproblems-problemphyxmini0273-timequantity}
  Read a duration in a selected time unit
\end{definition}
\begin{proof}
  This is the declared model definition of time Readout.
\end{proof}
\begin{definition}[angular Frequency Readout]
  \label{def:physics:phyx-mini-0273:phyxminiproblems-problemphyxmini0273-angularfrequencyreadout}
  \lean{PhyXMiniProblems.ProblemPhyXMini0273.angularFrequencyReadout}
  \uses{def:physics:phyx-mini-0273:phyxminiproblems-problemphyxmini0273-angularfrequencymagnitudequantity}
  Read an angular frequency in radians per selected time unit
\end{definition}
\begin{proof}
  This is the declared model definition of angular Frequency Readout.
\end{proof}
\begin{definition}[acceleration Magnitude Readout]
  \label{def:physics:phyx-mini-0273:phyxminiproblems-problemphyxmini0273-accelerationmagnitudereadout}
  \lean{PhyXMiniProblems.ProblemPhyXMini0273.accelerationMagnitudeReadout}
  \uses{def:physics:phyx-mini-0273:phyxminiproblems-problemphyxmini0273-accelerationmagnitudequantity}
  Read an acceleration magnitude in coherent selected length and time units
\end{definition}
\begin{proof}
  This is the declared model definition of acceleration Magnitude Readout.
\end{proof}
\begin{definition}[length In Meters]
  \label{def:physics:phyx-mini-0273:phyxminiproblems-problemphyxmini0273-lengthinmeters}
  \lean{PhyXMiniProblems.ProblemPhyXMini0273.lengthInMeters}
  \uses{def:physics:phyx-mini-0273:phyxminiproblems-problemphyxmini0273-lengthmagnitudequantity, def:physics:phyx-mini-0273:phyxminiproblems-problemphyxmini0273-lengthmagnitudereadout}
  Meter readout of a nonnegative length
\end{definition}
\begin{proof}
  This is the declared model definition of length In Meters.
\end{proof}
\begin{definition}[time In Seconds]
  \label{def:physics:phyx-mini-0273:phyxminiproblems-problemphyxmini0273-timeinseconds}
  \lean{PhyXMiniProblems.ProblemPhyXMini0273.timeInSeconds}
  \uses{def:physics:phyx-mini-0273:phyxminiproblems-problemphyxmini0273-timequantity, def:physics:phyx-mini-0273:phyxminiproblems-problemphyxmini0273-timereadout}
  Second readout of a duration
\end{definition}
\begin{proof}
  This is the declared model definition of time In Seconds.
\end{proof}
\begin{definition}[angular Frequency In Radians Per Second]
  \label{def:physics:phyx-mini-0273:phyxminiproblems-problemphyxmini0273-angularfrequencyinradianspersecond}
  \lean{PhyXMiniProblems.ProblemPhyXMini0273.angularFrequencyInRadiansPerSecond}
  \uses{def:physics:phyx-mini-0273:phyxminiproblems-problemphyxmini0273-angularfrequencymagnitudequantity, def:physics:phyx-mini-0273:phyxminiproblems-problemphyxmini0273-angularfrequencyreadout}
  SI radian-per-second readout of an angular frequency
\end{definition}
\begin{proof}
  This is the declared model definition of angular Frequency In Radians Per Second.
\end{proof}
\begin{definition}[acceleration In Meters Per Second Squared]
  \label{def:physics:phyx-mini-0273:phyxminiproblems-problemphyxmini0273-accelerationinmeterspersecondsquared}
  \lean{PhyXMiniProblems.ProblemPhyXMini0273.accelerationInMetersPerSecondSquared}
  \uses{def:physics:phyx-mini-0273:phyxminiproblems-problemphyxmini0273-accelerationmagnitudequantity, def:physics:phyx-mini-0273:phyxminiproblems-problemphyxmini0273-accelerationmagnitudereadout}
  SI meter-per-second-squared readout of an acceleration magnitude
\end{definition}
\begin{proof}
  This is the declared model definition of acceleration In Meters Per Second Squared.
\end{proof}
\begin{definition}[Graph Axis Variable]
  \label{def:physics:phyx-mini-0273:phyxminiproblems-problemphyxmini0273-graphaxisvariable}
  \lean{PhyXMiniProblems.ProblemPhyXMini0273.GraphAxisVariable}
  Variables printed beside the two graph axes
\end{definition}
\begin{proof}
  This is the declared model definition of Graph Axis Variable.
\end{proof}
\begin{definition}[Graph Axis Unit]
  \label{def:physics:phyx-mini-0273:phyxminiproblems-problemphyxmini0273-graphaxisunit}
  \lean{PhyXMiniProblems.ProblemPhyXMini0273.GraphAxisUnit}
  Unit abbreviations printed beside the two graph axes
\end{definition}
\begin{proof}
  This is the declared model definition of Graph Axis Unit.
\end{proof}
\begin{definition}[Marked Time Label]
  \label{def:physics:phyx-mini-0273:phyxminiproblems-problemphyxmini0273-markedtimelabel}
  \lean{PhyXMiniProblems.ProblemPhyXMini0273.MarkedTimeLabel}
  The special time symbol printed below the second crest
\end{definition}
\begin{proof}
  This is the declared model definition of Marked Time Label.
\end{proof}
\begin{definition}[Oscillator Kind]
  \label{def:physics:phyx-mini-0273:phyxminiproblems-problemphyxmini0273-oscillatorkind}
  \lean{PhyXMiniProblems.ProblemPhyXMini0273.OscillatorKind}
  The oscillator described in the problem statement
\end{definition}
\begin{proof}
  This is the declared model definition of Oscillator Kind.
\end{proof}
\begin{definition}[Circular Motion Kind]
  \label{def:physics:phyx-mini-0273:phyxminiproblems-problemphyxmini0273-circularmotionkind}
  \lean{PhyXMiniProblems.ProblemPhyXMini0273.CircularMotionKind}
  The auxiliary motion whose projection corresponds to the SHM
\end{definition}
\begin{proof}
  This is the declared model definition of Circular Motion Kind.
\end{proof}
\begin{definition}[Position Time Graph]
  \label{def:physics:phyx-mini-0273:phyxminiproblems-problemphyxmini0273-positiontimegraph}
  \lean{PhyXMiniProblems.ProblemPhyXMini0273.PositionTimeGraph}
  \uses{def:physics:phyx-mini-0273:phyxminiproblems-problemphyxmini0273-lengthmagnitudequantity, def:physics:phyx-mini-0273:phyxminiproblems-problemphyxmini0273-signedlengthquantity, def:physics:phyx-mini-0273:phyxminiproblems-problemphyxmini0273-timequantity, def:physics:phyx-mini-0273:phyxminiproblems-problemphyxmini0273-graphaxisvariable, def:physics:phyx-mini-0273:phyxminiproblems-problemphyxmini0273-graphaxisunit, def:physics:phyx-mini-0273:phyxminiproblems-problemphyxmini0273-markedtimelabel}
  The dimensionful data and qualitative labels present in the supplied graph. The position function is independent data; it is not defined as a cosine or as an answer formula
\end{definition}
\begin{proof}
  This is the declared model definition of Position Time Graph.
\end{proof}
\begin{definition}[SHMCircular Motion Setup]
  \label{def:physics:phyx-mini-0273:phyxminiproblems-problemphyxmini0273-shmcircularmotionsetup}
  \lean{PhyXMiniProblems.ProblemPhyXMini0273.SHMCircularMotionSetup}
  \uses{def:physics:phyx-mini-0273:phyxminiproblems-problemphyxmini0273-lengthmagnitudequantity, def:physics:phyx-mini-0273:phyxminiproblems-problemphyxmini0273-angularfrequencymagnitudequantity, def:physics:phyx-mini-0273:phyxminiproblems-problemphyxmini0273-accelerationmagnitudequantity, def:physics:phyx-mini-0273:phyxminiproblems-problemphyxmini0273-oscillatorkind, def:physics:phyx-mini-0273:phyxminiproblems-problemphyxmini0273-circularmotionkind, def:physics:phyx-mini-0273:phyxminiproblems-problemphyxmini0273-positiontimegraph}
  The spring-block SHM and the corresponding uniform circular motion. The two angular frequencies, the circle radius, and the radial acceleration are kept as independent physical quantities until related by the governing laws below
\end{definition}
\begin{proof}
  This is the declared model definition of SHMCircular Motion Setup.
\end{proof}
\begin{definition}[Matches Problem And Primary Figure]
  \label{def:physics:phyx-mini-0273:phyxminiproblems-problemphyxmini0273-matchesproblemandprimaryfigure}
  \lean{PhyXMiniProblems.ProblemPhyXMini0273.MatchesProblemAndPrimaryFigure}
  \uses{def:physics:phyx-mini-0273:phyxminiproblems-problemphyxmini0273-lengthmagnitudereadout, def:physics:phyx-mini-0273:phyxminiproblems-problemphyxmini0273-signedlengthreadout, def:physics:phyx-mini-0273:phyxminiproblems-problemphyxmini0273-timereadout, def:physics:phyx-mini-0273:phyxminiproblems-problemphyxmini0273-shmcircularmotionsetup}
  Problem-text and primary-image evidence. In the image the 8 cm and -8 cm numbers label grid lines; the curve reaches only 7 cm and -7 cm. The crest--trough--crest times establish a 40 ms crest-to-crest period
\end{definition}
\begin{proof}
  This is the declared model definition of Matches Problem And Primary Figure.
\end{proof}
\begin{definition}[Has Physical Parameters]
  \label{def:physics:phyx-mini-0273:phyxminiproblems-problemphyxmini0273-hasphysicalparameters}
  \lean{PhyXMiniProblems.ProblemPhyXMini0273.HasPhysicalParameters}
  \uses{def:physics:phyx-mini-0273:phyxminiproblems-problemphyxmini0273-lengthmagnitudereadout, def:physics:phyx-mini-0273:phyxminiproblems-problemphyxmini0273-timereadout, def:physics:phyx-mini-0273:phyxminiproblems-problemphyxmini0273-angularfrequencyreadout, def:physics:phyx-mini-0273:phyxminiproblems-problemphyxmini0273-shmcircularmotionsetup}
  Positivity and nondegeneracy conditions for the depicted motions
\end{definition}
\begin{proof}
  This is the declared model definition of Has Physical Parameters.
\end{proof}
\begin{definition}[Satisfies Simple Harmonic Motion Laws]
  \label{def:physics:phyx-mini-0273:phyxminiproblems-problemphyxmini0273-satisfiessimpleharmonicmotionlaws}
  \lean{PhyXMiniProblems.ProblemPhyXMini0273.SatisfiesSimpleHarmonicMotionLaws}
  \uses{def:physics:phyx-mini-0273:phyxminiproblems-problemphyxmini0273-timequantity, def:physics:phyx-mini-0273:phyxminiproblems-problemphyxmini0273-lengthmagnitudereadout, def:physics:phyx-mini-0273:phyxminiproblems-problemphyxmini0273-signedlengthreadout, def:physics:phyx-mini-0273:phyxminiproblems-problemphyxmini0273-timereadout, def:physics:phyx-mini-0273:phyxminiproblems-problemphyxmini0273-angularfrequencyreadout, def:physics:phyx-mini-0273:phyxminiproblems-problemphyxmini0273-shmcircularmotionsetup}
  The standard SHM laws with the phase chosen at the crest shown at t = 0: x(t) = A cos (omega t) and omega = 2 pi / T. These are general laws and contain neither the requested radial acceleration nor an answer choice
\end{definition}
\begin{proof}
  This is the declared model definition of Satisfies Simple Harmonic Motion Laws.
\end{proof}
\begin{definition}[Satisfies SHMUniform Circular Motion Correspondence]
  \label{def:physics:phyx-mini-0273:phyxminiproblems-problemphyxmini0273-satisfiesshmuniformcircularmotioncorrespondence}
  \lean{PhyXMiniProblems.ProblemPhyXMini0273.SatisfiesSHMUniformCircularMotionCorrespondence}
  \uses{def:physics:phyx-mini-0273:phyxminiproblems-problemphyxmini0273-shmcircularmotionsetup}
  The projection correspondence between SHM and uniform circular motion: the circle radius is the SHM amplitude, and both motions share angular frequency
\end{definition}
\begin{proof}
  This is the declared model definition of Satisfies SHMUniform Circular Motion Correspondence.
\end{proof}
\begin{definition}[Satisfies Uniform Circular Motion Radial Acceleration Law]
  \label{def:physics:phyx-mini-0273:phyxminiproblems-problemphyxmini0273-satisfiesuniformcircularmotionradialaccelerationlaw}
  \lean{PhyXMiniProblems.ProblemPhyXMini0273.SatisfiesUniformCircularMotionRadialAccelerationLaw}
  \uses{def:physics:phyx-mini-0273:phyxminiproblems-problemphyxmini0273-lengthmagnitudereadout, def:physics:phyx-mini-0273:phyxminiproblems-problemphyxmini0273-angularfrequencyreadout, def:physics:phyx-mini-0273:phyxminiproblems-problemphyxmini0273-accelerationmagnitudereadout, def:physics:phyx-mini-0273:phyxminiproblems-problemphyxmini0273-shmcircularmotionsetup}
  The centripetal-acceleration magnitude law a\_r = omega\textasciicircum{}2 r, stated in every coherent choice of length and time units. It is not specialized to the graph's numerical data or to choice C
\end{definition}
\begin{proof}
  This is the declared model definition of Satisfies Uniform Circular Motion Radial Acceleration Law.
\end{proof}
\begin{lemma}[radial Acceleration from graph amplitude and period]
  \label{lem:physics:phyx-mini-0273:phyxminiproblems-problemphyxmini0273-radialacceleration-from-graph-amplitude-and-period}
  \lean{PhyXMiniProblems.ProblemPhyXMini0273.radialAcceleration_from_graph_amplitude_and_period}
  \uses{def:physics:phyx-mini-0273:phyxminiproblems-problemphyxmini0273-lengthinmeters, def:physics:phyx-mini-0273:phyxminiproblems-problemphyxmini0273-timeinseconds, def:physics:phyx-mini-0273:phyxminiproblems-problemphyxmini0273-accelerationinmeterspersecondsquared, def:physics:phyx-mini-0273:phyxminiproblems-problemphyxmini0273-shmcircularmotionsetup, def:physics:phyx-mini-0273:phyxminiproblems-problemphyxmini0273-hasphysicalparameters, def:physics:phyx-mini-0273:phyxminiproblems-problemphyxmini0273-satisfiessimpleharmonicmotionlaws, def:physics:phyx-mini-0273:phyxminiproblems-problemphyxmini0273-satisfiesshmuniformcircularmotioncorrespondence, def:physics:phyx-mini-0273:phyxminiproblems-problemphyxmini0273-satisfiesuniformcircularmotionradialaccelerationlaw}
  Combining the three governing-law interfaces gives the radial acceleration in terms of the independently stored graph amplitude and period. This is a derived helper conclusion, not a premise field
\end{lemma}
\begin{proof}
  The conclusion follows from the governing relations and figure readouts named in the statement of radial Acceleration from graph amplitude and period.
\end{proof}
\begin{definition}[Answer Choice]
  \label{def:physics:phyx-mini-0273:phyxminiproblems-problemphyxmini0273-answerchoice}
  \lean{PhyXMiniProblems.ProblemPhyXMini0273.AnswerChoice}
  Labels printed beside the four candidate accelerations
\end{definition}
\begin{proof}
  This is the declared model definition of Answer Choice.
\end{proof}
\begin{definition}[meters Per Second Squared]
  \label{def:physics:phyx-mini-0273:phyxminiproblems-problemphyxmini0273-answerchoice-meterspersecondsquared}
  \lean{PhyXMiniProblems.ProblemPhyXMini0273.AnswerChoice.metersPerSecondSquared}
  \uses{def:physics:phyx-mini-0273:phyxminiproblems-problemphyxmini0273-answerchoice}
  Meter-per-second-squared value printed beside each answer label
\end{definition}
\begin{proof}
  This is the declared model definition of meters Per Second Squared.
\end{proof}
\begin{definition}[recorded Answer Choice]
  \label{def:physics:phyx-mini-0273:phyxminiproblems-problemphyxmini0273-recordedanswerchoice}
  \lean{PhyXMiniProblems.ProblemPhyXMini0273.recordedAnswerChoice}
  \uses{def:physics:phyx-mini-0273:phyxminiproblems-problemphyxmini0273-answerchoice}
  The answer label recorded in the supplied dataset
\end{definition}
\begin{proof}
  This is the declared model definition of recorded Answer Choice.
\end{proof}
\begin{definition}[Matches Displayed Answer]
  \label{def:physics:phyx-mini-0273:phyxminiproblems-problemphyxmini0273-matchesdisplayedanswer}
  \lean{PhyXMiniProblems.ProblemPhyXMini0273.MatchesDisplayedAnswer}
  \uses{def:physics:phyx-mini-0273:phyxminiproblems-problemphyxmini0273-accelerationmagnitudequantity, def:physics:phyx-mini-0273:phyxminiproblems-problemphyxmini0273-accelerationinmeterspersecondsquared, def:physics:phyx-mini-0273:phyxminiproblems-problemphyxmini0273-answerchoice}
  Agreement with an acceleration displayed to the nearest 100 m/s\textasciicircum{}2; the half-unit tolerance is 50 m/s\textasciicircum{}2 and applies uniformly to every choice
\end{definition}
\begin{proof}
  This is the declared model definition of Matches Displayed Answer.
\end{proof}
% --- Archon declaration topology end ---
% --- Archon physics formalization source end ---
